% Part: first-order-logic
% Chapter: natural-deduction
% Section: derivations

\documentclass[../../../include/open-logic-section]{subfiles}

\begin{document}

\iftag{FOL}
      {\olfileid{fol}{ntd}{der}}
      {\olfileid{pl}{ntd}{der}}

\olsection{\usetoken{P}{derivation}}

\begin{explain}
گفته‌ایم فرض چیست و قواعد استنتاج را نیز به دست داده‌ایم.
!!^{derivation}s در استنتاج طبیعی به‌طور استقرایی از این‌ها تولید
می‌شوند: هر !!{derivation} یا به‌تنهایی یک فرض است، یا از یک، دو یا سه
!!{derivation}s تشکیل می‌شود که استنتاجی درست پس از آن‌ها آمده است.
\end{explain}

\begin{defn}[!!^{derivation}]
\Article{derivation} \emph{!!{derivation}} از !!a{sentence}~$!A$ از
فرض‌های~$\Gamma$، درختی متناهی از !!{sentence}s است که شرایط زیر را
برآورده می‌کند:
\begin{enumerate}
\item بالاترین !!{sentence}s درخت یا عضو $\Gamma$ هستند یا به‌وسیلهٔ
  استنتاجی در درخت !!{discharged} می‌شوند.
\item پایین‌ترین !!{sentence} درخت~$!A$ است.
\item هر !!{sentence} درخت جز جملهٔ~$!A$ در پایین، مقدمهٔ کاربرد درستی
  از یک قاعدهٔ استنتاج است که نتیجه‌اش بی‌واسطه زیر آن !!{sentence} در
  درخت قرار دارد.
\end{enumerate}
آنگاه می‌گوییم $!A$~\emph{نتیجهٔ} !!{derivation} و $\Gamma$ فرض‌های
!!{undischarged} آن است.

اگر !!a{derivation} از $!A$ از~$\Gamma$ وجود داشته باشد، می‌گوییم
$!A$~\emph{!!{derivable}} از~$\Gamma$ است، یا به زبان نمادین:
$\Gamma \Proves !A$. اگر !!a{derivation} از~$!A$ وجود داشته باشد که در
آن هر فرض !!{discharged} شده است، می‌نویسیم~$\Proves !A$.
\end{defn}

\begin{ex}
هر فرض به‌تنهایی !!a{derivation} است. پس برای مثال، $!A$ به‌تنهایی
!!a{derivation} است و $!B$ به‌تنهایی نیز چنین است. می‌توانیم با
اعمال مثلاً قاعدهٔ $\Intro{\land}$ به این‌ها !!{derivation} تازه‌ای
به دست آوریم:
\begin{prooftree}
\AxiomC{$!A$}
\AxiomC{$!B$}
\RightLabel{\Intro{\land}}
\BinaryInfC{$!A \land !B$}
\end{prooftree}
قرار است این قواعد کلی باشند: می‌توانیم $!A$ و~$!B$ در آن را با هر
!!{sentence}sی، برای مثال با $!C$ و~$!D$، جایگزین کنیم. در این صورت
نتیجه $!C \land !D$ می‌بود و بنابراین
\begin{prooftree}
\AxiomC{$!C$}
\AxiomC{$!D$}
\RightLabel{\Intro{\land}}
\BinaryInfC{$!C \land !D$}
\end{prooftree}
یک !!{derivation} درست است. البته می‌توانیم فرض‌ها را نیز جابه‌جا کنیم
تا $!D$ نقش~$!A$ و $!C$ نقش~$!B$ را بازی کند. ازاین‌رو،
\begin{prooftree}
\AxiomC{$!D$}
\AxiomC{$!C$}
\RightLabel{\Intro{\land}}
\BinaryInfC{$!D \land !C$}
\end{prooftree}
نیز یک !!{derivation} درست است.

اکنون می‌توانیم قاعده‌ای دیگر، مثلاً $\Intro{\lif}$، را اعمال کنیم که
به ما اجازه می‌دهد یک شرطی را نتیجه بگیریم و هر فرضی را که با مقدم آن
شرطی یکسان است !!{discharge} کنیم. پس هر دو عبارت زیر !!{derivation}s
درستی خواهند بود:
\begin{prooftree}
\AxiomC{$\Discharge{!C}{1}$}
\AxiomC{$!D$}
\RightLabel{\Intro{\land}}
\BinaryInfC{$!C \land !D$}
\DischargeRule{\Intro{\lif}}{1}
\UnaryInfC{$!C \lif (!C \land !D)$}
\DisplayProof\bottomAlignProof
\AxiomC{$!C$}
\AxiomC{$\Discharge{!D}{1}$}
\RightLabel{\Intro{\land}}
\BinaryInfC{$!C \land !D$}
\DischargeRule{\Intro{\lif}}{1}
\UnaryInfC{$!D \lif (!C \land !D)$}
\end{prooftree}
آن‌ها به‌ترتیب نشان می‌دهند که $!D \Proves !C \lif (!C \land !D)$ و
$!C \Proves !D \lif (!C \land !D)$.

به یاد داشته باشید که رفع فرض‌ها اجازه است نه الزام: ناچار نیستیم فرض‌ها
را رفع کنیم. به‌ویژه، حتی اگر فرض‌ها در !!{derivation} حاضر نباشند نیز
می‌توانیم قاعده‌ای را اعمال کنیم. برای نمونه، عبارت زیر مجاز است، هرچند
هیچ فرض~$!A$ای برای !!{discharged} شدن وجود ندارد:
\begin{prooftree}
  \AxiomC{$!B$}
  \DischargeRule{\Intro{\lif}}{1}
  \UnaryInfC{$!A \lif !B$}
\end{prooftree}
\end{ex}

\end{document}
